\begin{problem}{Karesansui Garden}{standard input}{standard output}{2 seconds}{256 megabytes}

In the temple's karesansui (dry landscape) garden, $n$ gravel mounds stand in a row; mound $i$ is piled $a_i$ scoops high. For the evening meditation the head monk wants a stretch of at least $k$ consecutive mounds raked perfectly flat.

Two kinds of work are available, any number of times, in any order:

\begin{itemize}
  \item \textbf{Pay:} choose an index $i$ ($1 \le i \le n$) and a positive integer $c \le a_i$, pay the groundskeepers $c$ coins, and they carry off that many scoops: set $a_i := a_i - c$.
  \item \textbf{Rake:} choose three indices $l$, $m$, $r$ such that $1 \le l < m < r \le n$, drag the great garden rake from mound $l$ to mound $r$, and mound $m$ settles at the height of its taller support: set $a_m := \max(a_l, a_r)$. Raking costs nothing.
\end{itemize}

Find the minimum total number of coins the temple has to pay so that some contiguous segment of length at least $k$ of the row consists only of zeros.


\InputFile
The first line contains a single integer $t$ ($1 \le t \le 2 \cdot 10^5$) --- the number of test cases.

The first line of each test case contains two integers $n$ and $k$ ($1 \le k \le n \le 2 \cdot 10^5$) --- the size of the array and the minimum length of the required zero segment.

The second line of each test case contains $n$ integers $a_1, a_2, \ldots, a_n$ ($1 \le a_i \le 10^9$) --- the heights of the mounds.

It is guaranteed that the sum of $n$ over all test cases does not exceed $2 \cdot 10^5$.


\OutputFile
For each test case, print a single integer --- the minimum total cost needed to obtain a contiguous zero segment of length at least $k$.


\Example

\begin{example}
\exmpfile{example.01}{example.01.a}%
\end{example}

\Note
In the first test case, pay $1$ coin at mound $1$ ($a_1 := 1 - 1 = 0$) and $2$ coins at mound $3$ ($a_3 := 2 - 2 = 0$), then rake with $(l, m, r) = (1, 2, 3)$: mound $2$ settles at $\max(0, 0) = 0$ and the stretch $[1, 3]$ is flat, for a total of $3$ coins.

In the second test case, pay $10$ coins at mound $1$ and settle mound $3$ in two payments of $25$ and $5$ coins (so $a_3$ goes $30 \to 5 \to 0$); one rake with $(1, 2, 3)$ flattens mound $2$ for free --- $40$ coins in total.

In the third test case $k = 1$, so paying all $7$ coins for the only mound already suffices. In the fourth test case, paying mounds $2$ and $3$ ($1 + 1$ coins) flattens the stretch $[2, 3]$.


\end{problem}

